\section{Práctica: reproducir la esencia de la clase}

\subsection{Materiales y preparación}
Para reproducir esta clase primero hay que conseguir los recursos correspondientes: el video `audio.m4a`, `lecture09.srt`, la portada `cover.jpg` y cualquier slides/handout disponible. Después de limpiar los subtítulos a \texttt{subs\_clean.txt} (quitando líneas vacías y oraciones repetidas), se pueden localizar rápidamente los segmentos clave del docente según las marcas de tiempo. Los metadatos también deben quedar completos: autor, fecha de publicación, enlace del video y duración, para que la plantilla previa pueda citarlos directamente en el recuadro de portada. Si ya se dispone de las slides, se puede usar `magick -density 150 slide.pdf slide.jpg` para convertir cada página en imagen y colocarla en un `figure` como evidencia visual adicional; si no hay slides, hay que capturar los fotogramas clave del video y anotar el timecode correspondiente.

\begin{knowledgebox}{Preparación de slides y fotogramas clave}
\begin{itemize}
    \item Si las slides ya existen en PDF, priorizar la captura de las páginas importantes e insertarlas en un `figure`, anotando en el caption tanto el número de página como la indicación de tiempo.
    \item Si no hay slides, localizar en el video los momentos clave (por ejemplo, 00:22:10 o 00:35:40), capturar el fotograma, colocarlo en un `figure` y anotar el tiempo.
    \item Al citar las palabras originales con `\textit{}`, escribir la marca de tiempo en el formato `[HH:MM:SS]` para facilitar la referencia posterior.
\end{itemize}
\end{knowledgebox}

\subsection{Proceso de análisis y redacción}
El proceso de redacción de las notas se puede dividir en: leer los subtítulos limpios, elaborar el «Teaching Logic Outline», desarrollar cada sección según esa lógica, insertar recuadros y agregar figuras y resúmenes. Cada sección termina con un `\subsection{Resumen de la sección}`, lo que da lugar a una estructura reproducible.
\begin{itemize}
    \item Extraer palabras clave temáticas de \texttt{subs\_clean} y ordenarlas por temas como Datos, Algoritmo, Agente e Infraestructura.
    \item Citar las palabras originales del docente con `\textit{''}`, anotando la marca de tiempo e incorporándolas en la subsección correspondiente.
    \item Convertir los hallazgos clave en `importantbox`, `knowledgebox` y `warningbox`, para garantizar una jerarquía visual clara según el tipo de información.
    \item Por último, resumir toda la clase con una tabla de síntesis y lecturas adicionales, para que quien reproduzca el trabajo pueda revisarlo rápidamente.
\end{itemize}
\begin{knowledgebox}{Consejos de redacción al reproducir las notas}
\begin{itemize}
    \item Organizar las secciones según la lógica pedagógica, sin limitarse a acumular contenido en el orden temporal de los subtítulos.
    \item Cada idea importante debe repetirse, traducirse y citar su procedencia, conservando el sabor original con `\textit{}`.
    \item La evidencia visual debe acompañarse de una figura y una nota al pie con el tiempo; sirven la portada, las slides o los fotogramas clave.
    \item El `warningbox` se usa para advertir posibles malas interpretaciones, y el `importantbox` para resaltar conclusiones que representan un avance.
\end{itemize}
\end{knowledgebox}

\subsection{Evidencia visual y composición}
Esta sección no cuenta con slides, así que solo podemos usar fotogramas del video como fuente de imágenes: la portada `cover.jpg`, la presentación desde el lado del ponente y la interfaz de invocación de herramientas. Cada figura lleva una `\protect\footnotemark` que indica el intervalo de tiempo (por ejemplo, 00:00:05--00:00:13), y en el cuerpo del texto se explica por qué esa imagen es clave. Si en el futuro se encuentran las slides, se puede usar `magick -density ...` para convertir el PDF en `slide-*.jpg` e insertarlas en el `figure` en lugar del fotograma. Cada timecode debe quedar claramente escrito en el caption, para poder verificarlo después en una auditoría o revisión.

\begin{importantbox}{Requisitos de alta calidad para la evidencia visual}
Al capturar los fotogramas clave hay que mantener una nitidez de al menos 150 dpi, anotar el punto de tiempo en el caption y explicar el sentido de la imagen comparándola con las palabras originales del docente. Si se usan slides, hay que alinear el número de página con la línea de tiempo y evitar dejar textos de posición de enlace sin reemplazar o captions en blanco.
\end{importantbox}

\subsection{Proceso de reproducción y fotogramas clave}
El proceso de reproducción se puede dividir en tres etapas: 1) extraer los temas de \texttt{subs\_clean}, 2) distribuir los temas en las secciones según la lógica pedagógica, y 3) insertar imágenes con texto, recuadros y resúmenes. Cada sección debe registrar la marca de tiempo correspondiente, como `[00:40:12]` para la explicación de los pasos críticos del Agent Swarm. Si se detecta que el video tiene slides pero los subtítulos no las mencionan, aun así se pueden traducir manualmente y anotarlas en la descripción del `figure`, señalando «aún falta la explicación en inglés de la slide».

\begin{knowledgebox}{Fotogramas clave y documentación lista para auditoría}
\begin{itemize}
    \item Detrás de cada figura escribir el timecode, por ejemplo «intervalo de tiempo de la escena: 00:32:10--00:32:48».
    \item En la tabla resumen indicar los temas y la evidencia, como «Agent Swarm -> eficiencia 4.5 veces mayor -> figura [00:36:50]».
    \item Evitar textos de posición de metadatos sin reemplazar; toda la información de portada debe ser real.
\end{itemize}
\end{knowledgebox}

\subsection{Resumen de la sección}
Este proceso práctico: primero reunir los materiales y organizar los metadatos, después desglosar según la lógica pedagógica, y luego complementar cada capítulo con evidencia visual, citas y recuadros, permite producir notas K2.5 que cumplen con los criterios de calidad.
